\documentclass[11pt]{article}
\usepackage[margin=2.5cm]{geometry}
\usepackage{amsmath,amssymb}
\usepackage{booktabs}
\usepackage[hidelinks]{hyperref}

\newcommand{\het}{\nu}
\newcommand{\pd}{\partial}
\newcommand{\Op}{\operatorname{Op}^w}
\newcommand{\etal}{\eta_{\ell}}

\title{One-Way Helmholtz Propagation over Irregular Terrain\\ as a Weyl-Symbol Problem}
\author{Mikhail S. Lytaev\thanks{E-mail: mikelytaev@gmail.com. Code: \url{https://github.com/mikelytaev/wave-propagation}.}}
\date{}

\begin{document}
\maketitle

\begin{abstract}
A parabolic-equation solver for propagation over irregular terrain is built
from four ingredients: a one-way factorization of the Helmholtz operator, a
range step conjugated by the weights $B^{\pm1/4}$, a rational (Pad\'e)
propagator, and transparent boundaries. We develop all four in the calculus
of Weyl symbols for a propagation problem written in orthogonal
terrain-following coordinates. The
vertical operator $B$ of the Liouville-transformed equation
$\varphi_{\xi\xi}+B\varphi=0$ has an exact Weyl symbol whose principal part is
the local dispersion relation and which has no first-order part; the only
corrections are two zeroth-order terms of relative size $(\lambda/\ell)^2$.
The one-way factorization is the order-by-order solution of the operator
Riccati equation in the small parameter $\varepsilon=1/(k\ell)$; the
exponential propagator and its conjugation by the Weyl factors
$B^{\pm1/4}$ are statements of the
semiclassical functional calculus; the rational approximation is a pointwise function on phase
space whose error is budgeted on the transported support of the field; the
discrete scheme has a discrete symbol that coincides locally with that of the
flat scheme. The same calculus explains two obstructions that no grid
refinement removes: the shift map has a grid-independent error floor because
its one-way symbol has an odd part in the vertical wavenumber, and no local
top boundary can absorb grazing waves because every local condition is a
rational approximation of a square root up to its branch point at zero
vertical wavenumber.
\end{abstract}

%=====================================================================
\section{The Helmholtz equation}\label{sec:helm}

Consider the two-dimensional Helmholtz equation for a time-harmonic field
$\psi(x,z)$ ($e^{-i\omega t}$ dependence),
\begin{equation}\label{eq:helmholtz}
\nabla^2\psi+k^2(x,z)\psi=0 ,
\end{equation}
in the domain $\Omega=\{z>T(x)\}$ above the terrain profile $T(x)$, where $k=k_0n$
is the wavenumber and $n$ the modified refractive index including the Earth
flattening \cite{levy2000book}. The ground is either perfectly conducting,
with $\psi=0$ on $z=T(x)$ for horizontal (H) polarization and $\pd_n\psi=0$
for vertical (V) polarization, or it satisfies the Leontovich impedance
condition
\begin{equation}\label{eq:leontovich}
\pd_n\psi+\alpha\psi=0,\qquad \alpha=ik_0\Delta_g,\quad
\Delta_g=\begin{cases}\sqrt{\varepsilon_g-1}, & \text{H},\\
\sqrt{\varepsilon_g-1}/\varepsilon_g, & \text{V},\end{cases}
\end{equation}
with $\varepsilon_g$ the complex relative permittivity of the ground and the
normal directed into $\Omega$. A plane wave incident at the grazing angle
$\phi$ is reflected with $R(\phi)=(\sin\phi-\Delta_g)/(\sin\phi+\Delta_g)$.
Above some height the medium no longer depends on $x$, and the domain is
truncated there by a transparent boundary condition.

The numerical task is to march a one-way field along the terrain with a
wide-angle propagator of adjustable accuracy, imposing the ground condition
exactly and the top condition transparently. Two families of terrain
treatments dominate the practice \cite{beilis1979coupled,barrios1994tpem,donohue2000gtpe,holm2007wasm}.
The \emph{staircase} replaces the boundary by a sequence of steps; it imposes
the reflection coefficient of the wrong grazing angle and diffracts at every
corner. The \emph{shift map} $\zeta=z-T(x)$ straightens the boundary but moves
the whole column with the terrain, so the transformed medium depends on range
at every height, and the upper boundary condition can no longer be exact; its
wide-angle versions carry a slope error that survives grid refinement
\cite{wang2020smpe,zhou2017lfhighangle,gilbert2022wabt}. Conformal maps are
nowhere the identity and share the same defect.

The alternative developed in \cite{lytaev2020awpl,lytaev2018awpl} and here is
a map that is the identity above a transition layer, so that the exterior
stays range-independent and the discrete nonlocal boundary condition (NLBC)
of \cite{lytaev2018awpl} remains exact; below the layer, the map keeps the
structure of the flat problem on which the rational propagator relies. This
paper derives the whole scheme in one language: the calculus of Weyl symbols.
Every ingredient of the method---the terrain-following coordinates, the
one-way factorization, the weighted range step, the rational propagator and
its error budget, the discrete scheme, the dispersion analysis, and the two
fundamental obstructions (shift-map floor, grazing-wave absorption)---becomes
a statement about a function on phase space. The calculus also ranks the
terms the method neglects and states precisely when the neglect is justified.

%=====================================================================
\section{Orthogonal terrain-following coordinates}\label{sec:map}

\subsection{Level curves}
Let $(x,z)=(X(\xi,\eta),Z(\xi,\eta))$ map the computational rectangle onto
$\Omega$, with the ground at $\eta=0$. The coordinate lines $\eta=\text{const}$
are the level curves
\begin{equation}\label{eq:family}
z=F(x,\eta)=\eta+z_{\rm ref}+u(x,\eta),\qquad
\hat u(\kappa,\eta)=\hat b(\kappa)\,w(\eta/L)\,e^{-|\kappa|\eta},
\end{equation}
where $b=T-z_{\rm ref}$, hats denote the Fourier transform in $x$,
$w(t)=1-t^3(10-15t+6t^2)$ on $0\le t\le1$ and $w=0$ above, and $L$ is the
height of the transition layer. The construction has four properties:
$F(x,0)=T(x)$ exactly; $F=\eta+z_{\rm ref}$ for $\eta\ge\etal=L$, so the
coordinates are Cartesian above the layer; $u$ is harmonic near the ground, so
relief of horizontal scale $\ell$ decays over a height of about $\ell$; and
the map is regular whenever $F_\eta>0$, which bounds the smoothed slopes. The
terrain samples are pre-filtered with a Gaussian of width $\sigma$, so the
metric entering the analysis below is smooth.

\subsection{Trajectories and metrics}
The lines $\xi=\text{const}$ are the orthogonal trajectories of the level
curves. Writing the tangents as $\mathbf a_\xi=P(1,F_x)$ and
$\mathbf a_\eta=(X_\eta,F_xX_\eta+F_\eta)$ with $P=X_\xi$, orthogonality
$\mathbf a_\xi\cdot\mathbf a_\eta=0$ gives an initial-value problem for each
trajectory,
\begin{equation}\label{eq:traj}
X_\eta=G(X,\eta)\equiv-\frac{F_xF_\eta}{1+F_x^2},\qquad
X(\xi,\etal)=\xi ,
\end{equation}
integrated downward from the top of the layer, where the lines are vertical
and are labeled by their abscissa. The scale factors are the lengths of the
tangents,
\begin{equation}\label{eq:metric}
h_\eta=\frac{F_\eta}{\sqrt{1+F_x^2}},\qquad
h_\xi=P\sqrt{1+F_x^2},\qquad
P_\eta=G_x(X,\eta)\,P ,
\end{equation}
with $P(\xi,\etal)=1$: $h_\eta$ is the vertical spacing $F_\eta\,d\eta$
projected on the normal tilted by $\arctan F_x$, and $h_\xi$ converts the
trajectory spacing $P\,d\xi$ to arc length along the level curve. Their ratio
$r=h_\eta/h_\xi$ is the anisotropy of the computational cell; $r=1$ for a
conformal map. The arc length is $ds^2=h_\xi^2d\xi^2+h_\eta^2d\eta^2$.

\subsection{The equation in the coordinates}
The divergence theorem turns \eqref{eq:helmholtz} into
\begin{equation}\label{eq:helm_orth}
\pd_\xi\bigl(r\,\pd_\xi\psi\bigr)+\pd_\eta\!\left(\frac1r\,\pd_\eta\psi\right)
+h_\xi h_\eta k^2\psi=0 .
\end{equation}
On $\eta=0$ the normal derivative is $\pd_n=h_\eta^{-1}\pd_\eta$, so the
impedance condition \eqref{eq:leontovich} becomes the Robin condition
$\pd_\eta\psi+\gamma\psi=0$ with $\gamma(\xi)=\alpha\,h_\eta(\xi,0)$; Dirichlet
and Neumann conditions are invariant. Above $\etal$ both scale factors are $1$
and the equation is the flat one, which is what makes the exterior
range-independent and the NLBC exact. Between the ground and $\etal$, the
coefficients $r$, $h_\xi$, $h_\eta$ vary on the smoothing length $\ell$ and on
the transition height $L$; the rest of the paper treats that variation.

%=====================================================================
\section{The Weyl symbol of $B$}\label{sec:symbol}

\subsection{Liouville transformation}
The range dependence of $r$ enters \eqref{eq:helm_orth} through
$\pd_\xi(r\pd_\xi)$. Setting $\psi=r^{-1/2}\varphi$ removes it:
\begin{equation}\label{eq:liou}
\varphi_{\xi\xi}+B\varphi=0,\qquad
B\varphi=r^{-1/2}\pd_\eta\!\left(r^{-1}\pd_\eta(r^{-1/2}\varphi)\right)
+(k^2h_\xi^2-V)\varphi ,
\end{equation}
with $V=(r^{1/2})_{\xi\xi}/r^{1/2}$, a term of relative size
$(\lambda/\ell)^2$. The operator $B$ acts in $\eta$ with $\xi$ as a parameter
and is symmetric in $L^2(d\eta)$.

\subsection{Quantization}
Let $D=-i\pd_\eta$ and let $q$ be the variable dual to $\eta$; $q$ is the
computational vertical wavenumber. The Weyl quantization of a symbol
$a(\eta,q)$ is
\[
\Op(a)u(\eta)=\frac1{2\pi}\iint e^{i(\eta-\eta')q}\,
a\!\left(\tfrac{\eta+\eta'}2,q\right)u(\eta')\,d\eta'\,dq .
\]
Real symbols give symmetric operators. Composition is the Moyal product
\begin{equation}\label{eq:moyal}
a\#b=ab+\frac i2\{a,b\}
-\frac18\bigl(a_{qq}b_{\eta\eta}-2a_{q\eta}b_{q\eta}+a_{\eta\eta}b_{qq}\bigr)+\dots,
\qquad \{a,b\}=a_\eta b_q-a_qb_\eta ,
\end{equation}
whose odd-order terms are antisymmetric in $(a,b)$. Hence $a\#a=a^2+O(\pd^4)$
and the symbol of a commutator is $i\{a,b\}$ to leading order.

\paragraph{The small parameter.} Let $\ell$ be the shortest length on which
the metric varies: the smoothing length of the terrain, or the transition
height $L$. The expansion parameter is
\begin{equation}\label{eq:eps}
\varepsilon=\frac1{k\ell} .
\end{equation}
In rescaled variables $\eta/\ell$, $q/k$ every derivative of a coefficient
brings one factor $\varepsilon$; this is the semiclassical calculus with
$\hbar=\varepsilon$ \cite{zworski,dimassi}. Symbols are graded
$a=a_0+\varepsilon a_1+\varepsilon^2a_2+\dots$, with the physical variables
kept in the formulas.

\subsection{The exact symbol}
For any smooth $g$, the identities
$\Op(g\,q^2)=\tfrac14(D^2g+2DgD+gD^2)$ and
$D^2g+gD^2=2DgD-g''$ give
\begin{equation}\label{eq:DgD}
DgD=\Op\!\left(g\,q^2+\tfrac14 g''\right).
\end{equation}
With $f=r^{-1/2}$ and $g=r^{-1}$, moving $f$ through $D$ gives
$fDgDf=D(f^2g)D-(ff'g)'+gf'^2$. The derivative part of $B$ \eqref{eq:liou} is
$-fDgDf$; applying \eqref{eq:DgD} to $f^2g=r^{-2}$, the second derivatives of
$r$ cancel and only a square of $r_\eta$ survives:
\begin{equation}\label{eq:symbol}
\boxed{\;b(\xi,\eta,q)=k^2h_\xi^2-\frac{q^2}{r^2}-V-\frac{r_\eta^2}{4r^4}\;},
\qquad B=\Op(b)\ \text{exactly}.
\end{equation}
The identity behind \eqref{eq:symbol},
$r^{-1/2}(r^{-1}(r^{-1/2}u)')'=(r^{-2}u')'+\tfrac14(r^{-2})''u
-\tfrac14r'^2r^{-4}u$, is elementary but was checked numerically to working
precision.

\paragraph{Principal symbol.} The principal part
\begin{equation}\label{eq:b0}
b_0=k^2h_\xi^2-\frac{q^2}{r^2}
\end{equation}
is the local dispersion relation: with $\kappa$ dual to $\xi$, the
characteristic set of \eqref{eq:liou} is $\kappa^2=b_0$, i.e.\
$(\kappa/h_\xi)^2+(q/h_\eta)^2=k^2$, the local plane-wave dispersion. With
$q=kh_\eta\sin\theta$, $\theta$ the angle to the level curve,
\begin{equation}\label{eq:angle}
b_0=k^2h_\xi^2\cos^2\theta,\qquad \sqrt{b_0}=kh_\xi\cos\theta .
\end{equation}

\paragraph{No first-order part.} The symbol \eqref{eq:symbol} is real and has
nothing of order $\varepsilon^1$: the symbol form of ``$B$ is symmetric''. It
matters because a subprincipal (imaginary, order-$\varepsilon q$) term would
transport amplitude across $\eta$ and make $B$ non-normal. The original
equation \eqref{eq:helm_orth} divided by $r$ does have first-order terms in
both variables, $(r_\xi/r)\pd_\xi$ and $\pd_\eta(r^{-1})\,\pd_\eta/r^{-1}$;
the Liouville substitution is the gauge that removes them. The price is two
zeroth-order terms of order $\varepsilon^2$:
\begin{itemize}
\item $V=(r^{1/2})_{\xi\xi}/r^{1/2}$, from the $\xi$ direction. Since $\xi$ is
  the evolution variable and is not quantized, it is written explicitly and
  computed by a second difference along the columns.
\item $W=r_\eta^2/(4r^4)$, from the $\eta$ direction. It is never assembled
  separately: the scheme discretizes the operator form
  $r\pd_\eta r^{-1}\pd_\eta$, which contains $W$ exactly.
\end{itemize}
Both are $O((\lambda/\ell)^2)$ relative to $k^2h_\xi^2$.

\paragraph{Stretching is a metric.} A vertical stretching $\eta=Z(\eta')$
with $s=dZ/d\eta'$ multiplies $h_\eta$ by $s$, i.e.\ replaces $r\to rs$ in
\eqref{eq:symbol}; then $q'/(h_\eta s)$ is the physical vertical wavenumber.
Absorbing layers and graded grids are therefore metric operations and stay
inside the calculus.

%=====================================================================
\section{One-way factorization and the $\varepsilon$-factorization}\label{sec:factor}

\subsection{The Riccati equation}
Seek a one-way equation $\varphi_\xi=i\Lambda\varphi$ with
$\Lambda=\Op(\lambda)$. Substituting into \eqref{eq:liou} gives the operator
Riccati equation
\begin{equation}\label{eq:riccati}
\Lambda^2-i\,\pd_\xi\Lambda=B
\quad\Longleftrightarrow\quad
\lambda\#\lambda-i\,\pd_\xi\lambda=b .
\end{equation}
If $B$ does not depend on $\xi$, the exact solution is the spectral square
root $\Lambda=B^{1/2}$. If it does, no exact solution representable as a
function of $B$ exists; the best one can do is solve \eqref{eq:riccati}
asymptotically.

\paragraph{The $\varepsilon$-factorization.}
The \emph{$\varepsilon$-factorization} is the asymptotic solution of
\eqref{eq:riccati}; in a range-dependent medium it is the working definition
of ``the'' one-way operator. It rests on one assumption about the medium and
three explicit rules; everything else is computation.

\emph{Assumption (slow variation).} The coefficients entering
\eqref{eq:symbol}---$h_\xi$, the anisotropy $r=h_\eta/h_\xi$, and $k$---vary
in \emph{both} $\xi$ and $\eta$ on scales no shorter than $\ell$; $\ell$ is
the shortest of these scales, the smoothing length of the terrain or the
transition height $L$. In the rescaled variables $x=\eta/\ell$, $p=q/k$
every symbol below is then a smooth function of $(\xi/\ell,x,p)$ with
bounded derivatives, the small parameter is
$\varepsilon=1/(k\ell)=\lambda/\ell$, and on the band of interest the
principal symbol stays away from its turning set $b_0=0$.

\emph{Rules.}
(i)~The unknown is expanded as
\begin{equation}\label{eq:expansion}
\lambda=\lambda_0+\lambda_1+\lambda_2+\dots,
\end{equation}
where $\lambda_j$ is of order $\varepsilon^j$: relative to $\lambda_0$,
each $\eta$-derivative or $\xi$-derivative multiplies its size by another
$\varepsilon$, and a $q$-derivative leaves it unchanged.
(ii)~The Moyal product is graded by the same count: in the rescaled
variables,
$a\#b=ab+\tfrac{i\varepsilon}2\{a,b\}+O(\varepsilon^2)$, the term with $m$
transverse derivatives being of order $\varepsilon^m$ relative to $ab$ (in
the physical variables of \eqref{eq:moyal} the powers of $\varepsilon$ are
implicit in the derivatives).
(iii)~Each $\pd_\xi$ of an order-$\varepsilon^j$ symbol is placed in order
$\varepsilon^{j+1}$: by the assumption a range derivative of a coefficient
is of size $1/\ell=\varepsilon k$, so
$\pd_\xi\lambda_j\sim\varepsilon\,k\,\lambda_j$.

\emph{Procedure.} Substitute \eqref{eq:expansion} into
\eqref{eq:riccati}, expand each product by rule~(ii) and each derivative by
rule~(iii), and equate the coefficients of equal powers of $\varepsilon$.
At order $\varepsilon^0$ the only unknown is $\lambda_0$. At each order
$\varepsilon^j$, $j\ge1$, the only term of top order containing the new
unknown is $2\lambda_0\lambda_j$; every other term is built from products,
brackets and $\xi$-derivatives of the known $\lambda_0,\dots,\lambda_{j-1}$
and of $b$. Since $\lambda_0\neq0$ on the band, each equation determines
$\lambda_j$ uniquely, and the recursion closes: the truncated sum
$\lambda^{(J)}=\lambda_0+\dots+\lambda_J$ satisfies \eqref{eq:riccati} up
to $O(\varepsilon^{J+1})$---this is what ``the field is factorized into
forward and backward waves'' means, to that order. The accuracy is set by
$\varepsilon=\lambda/\ell$, the smoothness of the medium on the wavelength
scale, \emph{not} by any grid. The next three paragraphs carry this out to
$j=2$.

\paragraph{Order 0: the principal symbol.} First grade the symbol
\eqref{eq:symbol} itself:
\[
b=b_0+b_2,\qquad
b_0=k^2h_\xi^2-\frac{q^2}{r^2},\qquad
b_2=-V-\frac{r_\eta^2}{4r^4},
\]
with $b_0$ of order $0$ and $b_2$ of order $\varepsilon^2$. The count is the
one of Section~\ref{sec:symbol}: each $\eta$-derivative of a metric
coefficient carries one factor $\varepsilon$; $V$ and $W=r_\eta^2/(4r^4)$
contain one and two of them and are of order $\varepsilon^2$, and
\eqref{eq:symbol} has no order-$\varepsilon$ term at all. At order
$\varepsilon^0$ the Riccati equation \eqref{eq:riccati} reads
$\lambda_0\#\lambda_0=b_0$, and
\[
\lambda_0\#\lambda_0=\lambda_0^2+O(\varepsilon^2):
\]
the first Moyal bracket vanishes identically,
$\{\lambda_0,\lambda_0\}=0$, and every summand of the second Moyal term
$\lambda_0\#\lambda_0\bigr|_{2}
=-\tfrac18\bigl[(\lambda_0)_{qq}(\lambda_0)_{\eta\eta}
-2(\lambda_0)_{q\eta}^2+(\lambda_0)_{\eta\eta}(\lambda_0)_{qq}\bigr]$
contains exactly two $\eta$-derivatives, hence is of order $\varepsilon^2$
(it enters the next order). The order-$\varepsilon^0$ equation is therefore
the algebraic equation
\[
\lambda_0^2=b_0 ,
\]
with exactly two roots, $\pm\sqrt{b_0}$; separating them is what the
factorization means, and the separation is the radiation condition. Fix the
\emph{outgoing} root
\[
\lambda_0=\sqrt{b_0},\qquad \operatorname{Im}\lambda_0\ge0 ,
\]
i.e.\ $\lambda_0=+|b_0|^{1/2}$ where $b_0\ge0$ and
$\lambda_0=i\,|b_0|^{1/2}$ where $b_0<0$. On the propagating part of the
band, $b_0\ge k^2h_\xi^2\cos^2\theta_{\max}>0$ for angles
$\theta\le\theta_{\max}$ within the aperture, $\lambda_0$ is real, positive
and a smooth function of $(\eta,q)$, as $b_0$ is bounded away from zero
there; on the evanescent part, $b_0<0$, one has
$\operatorname{Im}\lambda_0=|b_0|^{1/2}>0$, so the component
$e^{i\xi\lambda_0}$ carries the factor $e^{-\xi\operatorname{Im}\lambda_0}$
and decays as $\xi$ grows. The single statement
$\operatorname{Im}\lambda_0\ge0$ covers both: forward evolution damps
evanescent components and leaves propagating ones oscillatory. The root is
smooth on the whole band as long as the turning set $b_0=0$ stays outside
it; at $b_0=0$ the two roots coalesce, $\sqrt{b_0}$ loses smoothness, and no
branch with $\operatorname{Im}\lambda_0\ge0$ continues through---the
turning-point limit discussed at the end of this section. Note also that
$\lambda_0\neq0$ everywhere on the band, which makes every division by
$\lambda_0$ below legitimate.

\paragraph{Order 1: the transport term.} Insert
$\lambda=\lambda_0+\lambda_1+O(\varepsilon^2)$ into the Moyal product and
keep the terms through order $\varepsilon$. Three counts are needed.
First, $\lambda_0\#\lambda_0=\lambda_0^2+O(\varepsilon^2)$, by the
previous paragraph. Second,
$\lambda_0\#\lambda_1+\lambda_1\#\lambda_0
=2\lambda_0\lambda_1+O(\varepsilon^2)$: the first Moyal brackets of the two
products enter with opposite signs and cancel identically,
$\tfrac i2\{\lambda_0,\lambda_1\}+\tfrac i2\{\lambda_1,\lambda_0\}=0$, and
the second Moyal terms of both products are $O(\varepsilon^3)$---each of
their three summands distributes exactly two $\eta$-derivatives between the
factors, and $\lambda_1$ is itself of order $\varepsilon$, e.g.\
$(\lambda_0)_{qq}(\lambda_1)_{\eta\eta}=O(1)\cdot O(\varepsilon^3)$,
$(\lambda_0)_{q\eta}(\lambda_1)_{q\eta}
=O(\varepsilon)\cdot O(\varepsilon^2)$,
$(\lambda_0)_{\eta\eta}(\lambda_1)_{qq}
=O(\varepsilon^2)\cdot O(\varepsilon)$, a $q$-derivative carrying no
$\varepsilon$. Third, $\pd_\xi\lambda=\pd_\xi\lambda_0+O(\varepsilon^2)$,
because
$\lambda_0$ depends on $\xi$ only through the metric coefficients, which
vary on the scale $\ell$, so that $\pd_\xi\lambda_0
=\pd_\xi b_0/(2\lambda_0)$ is of order $\varepsilon$. The
order-$\varepsilon^1$ part of \eqref{eq:riccati} is consequently
\[
2\lambda_0\lambda_1-i\,\pd_\xi\lambda_0=0 ,
\]
and since $\lambda_0\neq0$ on the band, it has exactly one solution:
\begin{equation}\label{eq:lambda1}
\lambda_1=\frac i2\,\frac{\pd_\xi\lambda_0}{\lambda_0}
=\frac i4\,\frac{\pd_\xi b_0}{b_0}
=\frac i4\,\frac{\pd_\xi b}{b} ,
\end{equation}
where the last equality holds at the accuracy of the expansion: the
difference is $(i/4)\,\pd_\xi b_2/b_0$, of order $\varepsilon^3$, below the
resolution of an order-$\varepsilon$ symbol. Three consequences.
(i)~\emph{Uniqueness.} Once the branch of $\lambda_0$ is fixed, every
$\lambda_j$ for $j\ge1$ is determined recursively by a linear equation
$2\lambda_0\lambda_j=(\text{known symbols of order}\le j)$ with
$\lambda_0\neq0$: the $\varepsilon$-factorization contains no freedom
beyond the radiation condition. (ii)~\emph{Ordering independence.}
$\lambda_1$ is the order-$\varepsilon$ symbol of the classical factor
$\tfrac i4B^{-1}B_\xi$, in the precise sense that
\[
\operatorname{symbol}\Bigl(\tfrac i4B^{-1}B_\xi\Bigr)
=\tfrac i4\,b^{-1}\# b_\xi
=\tfrac i4\Bigl[\frac{b_\xi}{b}
+\frac i2\,\{b^{-1},b_\xi\}\Bigr]+O(\varepsilon^2)
=\lambda_1+O(\varepsilon^2),
\]
because the bracket $\{b^{-1},b_\xi\}$ contains one $\eta$-derivative of
the order-$\varepsilon$ symbol $b_\xi$ and is itself of order
$\varepsilon^2$, while the second Moyal term is of order $\varepsilon^3$.
So ``$B$ and $B_\xi$ treated as commuting'' is exact at this order: every
operator ordering of the product $B^{-1}B_\xi$ has the same
order-$\varepsilon$ symbol, and ordering enters only the
order-$\varepsilon^2$ terms of \eqref{eq:lambda2} and the commutator
$i\{b,b_\xi\}$. (iii)~\emph{Phase structure.} $\lambda_1$ is purely
imaginary---$i$ times the real symbol $\pd_\xi b/(4b)$---so it rescales the
amplitude of $e^{i\xi\,\Op(\lambda)}$ and does not contribute to its phase;
it is precisely the term eliminated by the change of variables of
Section~\ref{sec:propagator}.

\paragraph{Order 2.} Collecting the $\varepsilon^2$ terms of
\eqref{eq:riccati}, including the second Moyal term of $\lambda_0\#\lambda_0$,
\begin{equation}\label{eq:lambda2}
\lambda_2=\frac1{2\sqrt b}\left[
\frac{5\,b_\xi^2}{16\,b^2}-\frac{b_{\xi\xi}}{4b}
+\frac14\Bigl(\pd_q^2\lambda_0\,\pd_\eta^2\lambda_0
-(\pd_q\pd_\eta\lambda_0)^2\Bigr)\right].
\end{equation}
Three remarks. First, if $b$ does not depend on $\eta$ the bracket vanishes
and what remains is the classical second-order WKB correction
$\tfrac5{32}Q'^2Q^{-5/2}-\tfrac18Q''Q^{-3/2}$ for $y''+Qy=0$: a check of the
expansion. Second, the bracket is the phase-space curvature of the one-way
symbol; together with the ordering ambiguity of $\tfrac i4B^{-1}B_\xi$
(a term $\propto\{b^{-1},b_\xi\}$) it is precisely the ``second-order terms
and commutators $[B,B_\xi]$'' of the usual formal expansions, whose symbol is
$i\{b,b_\xi\}$. Third, all of it is $O(\varepsilon^2)$; over a transition
region of length $\ell_\xi$ these terms accumulate a phase of order
$k\ell_\xi\varepsilon^2$, which for smooth terrain is tiny.

\paragraph{Limits.} The expansion needs $\varepsilon\ll1$---terrain smoothing
and the transition layer are what keep the metric smooth on the wavelength
scale---and it fails where $b_0=0$, i.e.\ $\theta\to90^\circ$ to the marching
direction, because $\sqrt{b_0}$ is not smooth there. Within an aperture
$\theta\le\theta_{\max}$ one has
$b_0\ge k^2h_\xi^2\cos^2\theta_{\max}>0$, keeping the field away from the
turning set. What the factorization drops quantitatively---a backward wave of
size $O(\varepsilon)$ locally but $O(\varepsilon^\infty)$ genuinely
scattered---is measured in Section~\ref{sec:backscatter}.

%=====================================================================
\section{The exponential propagator in Weyl symbols}\label{sec:propagator}

\subsection{Functional calculus}
For $B=\Op(b)$ symmetric and a smooth function $F$, the semiclassical
functional calculus gives
\begin{equation}\label{eq:funcalc}
F(B)=\Op\!\left(F\circ b\right)+O(\varepsilon^2)
\end{equation}
\cite{helffer,dimassi}; in the Weyl calculus the $O(\varepsilon)$ correction
vanishes. Applied to the one-step propagator with $L=B/\beta^2-1$ and a
reference wavenumber $\beta$,
\begin{equation}\label{eq:expsymbol}
\operatorname{symbol}\bigl(e^{i\beta\Delta\xi(\sqrt{1+L}-1)}\bigr)
=e^{i\Delta\xi(\sqrt b-\beta)}+O(\varepsilon^2):
\end{equation}
\emph{the exponential propagator is, to leading order, multiplication by the
ordinary function $e^{i\Delta\xi(\sqrt b-\beta)}$ on phase space.} All errors
introduced by treating a variable-coefficient step as a frozen one live in
this remainder.

\subsection{The transport equation and the change of variables}
Everything in this subsection is a computation on symbols, via the Moyal
product \eqref{eq:moyal}; no operator algebra is used beyond the calculus
\eqref{eq:funcalc} itself.

\paragraph{The transport equation.} Take the factorization of Section
\ref{sec:factor} up to $\varepsilon^2$,
$\lambda^{(2)}=\lambda_0+\lambda_1+\lambda_2$, with
$\lambda_0=\sqrt b$, $\lambda_1=i\rho$ by \eqref{eq:lambda1},
$\rho=\pd_\xi b/(4b)$, and $\lambda_2$ as in \eqref{eq:lambda2}. Then
\eqref{eq:liou} is equivalent to the first-order transport equation
\begin{equation}\label{eq:transport}
\varphi_\xi=\Op\bigl(i\lambda^{(2)}\bigr)\varphi+O(\varepsilon^3)
=i\,\Op(\sqrt b)\,\varphi-\Op(\rho)\,\varphi+i\,\Op(\lambda_2)\,\varphi
+O(\varepsilon^3):
\end{equation}
a frozen-medium oscillation, a real amplitude term of order $\varepsilon$,
and an oscillatory remainder of order $\varepsilon^2$. (In operator notation
this is the equation
$\varphi_\xi=i\sqrt B\varphi-\tfrac14B^{-1}B_\xi\varphi+\dots$, the symbol
of $\tfrac14B^{-1}B_\xi$ being $\rho+O(\varepsilon^2)$.) The term
$-\Op(\rho)$ is what a step must reproduce; the term $i\Op(\lambda_2)$ is
the remainder quantified in Section~\ref{sec:backscatter}.

\paragraph{Change of variables.} Seek $\chi=\Op(w)\varphi$, with an
order-$0$ symbol $w$ to be determined, such that $\chi$ satisfies the
frozen-medium transport
$\chi_\xi=i\Op(\lambda_0)\chi+O(\varepsilon^2)$. Differentiating and using
\eqref{eq:transport} at order $\varepsilon^1$,
\[
\chi_\xi=\Op\bigl(w_\xi+w\#(i\lambda_0-\rho)\bigr)\varphi
+O(\varepsilon^2),\qquad
i\Op(\lambda_0)\chi=\Op(i\lambda_0\#w)\varphi,
\]
so the requirement reduces to the vanishing of the mismatch
\[
w_\xi+w\#(i\lambda_0-\rho)-i\lambda_0\#w
=w_\xi-\rho\,w-\{w,\lambda_0\}+O(\varepsilon^2).
\]
The change of variables is therefore governed by a linear differential
equation on phase space,
\begin{equation}\label{eq:weq}
\pd_\xi w=\rho\,w+\{w,\lambda_0\}+O(\varepsilon^2):
\end{equation}
a scalar growth $\rho$ and a transport along the Hamiltonian flow of
$\lambda_0$, the bicharacteristic field \eqref{eq:rays}. Its general
solution is $w=F(b)\,e^{\int\rho\,d\xi}+O(\varepsilon^2)$: any $F(b)$
annihilates the bracket exactly, $\{F(b),\lambda_0\}=F'\{\,b,\sqrt b\,\}
=F'\lambda_0'\{b,b\}=0$, two functions of one symbol commuting in the first
Moyal order, and the remaining factor obeys the scalar equation
$w_\xi=\rho w$. The free $F(b)$ is itself a change of variables commuting
with $i\Op(\lambda_0)$; fixing it by its value at the initial column
($F\equiv1$) gives the canonical choice
\[
w=b^{1/4},\qquad\text{i.e.}\qquad \chi=\Op\bigl(b^{1/4}\bigr)\varphi ,
\]
and this $w$ solves \eqref{eq:weq} identically:
$\{b^{1/4},\sqrt b\}=0$ and
$\pd_\xi b^{1/4}=\tfrac14b^{-3/4}\pd_\xi b=\rho\,b^{1/4}$. The change of
variables therefore introduces no error of its own at this order.

\paragraph{Transformed equation.} In $\chi$,
\[
\chi_\xi=i\Op(\lambda_0)\,\chi+O(\varepsilon^2):
\]
the same equation as in a frozen medium. The order-$\varepsilon$ term
\eqref{eq:lambda1} is removed, not approximated. Since $\lambda_0$ is real,
$\Op(\lambda_0)$ is symmetric and the generator $i\Op(\lambda_0)$ is
skew-adjoint, so $\tfrac d{d\xi}\|\chi\|^2=0$ up to $O(\varepsilon^2)$; for
$\varphi$ itself the trace identity
$\langle u,\Op(a)u\rangle=\iint a\,n_u\,d\eta\,dq$, with $n_u$ the Wigner
transform of $u$, gives
$\tfrac d{d\xi}\|\varphi\|^2=-2\iint\rho\,n_\varphi+O(\varepsilon^2)$:
the weight carries the amplitude. The order-$\varepsilon^2$ remainder
$i\Op(\lambda_2)$ is not removed at this order; what it contains is the
subject of Section~\ref{sec:backscatter}.

\paragraph{The one-step propagator.} Advance $\chi$ over
$[\xi_n,\xi_{n+1}]$ by \eqref{eq:expsymbol} with the symbol frozen at the
midpoint column---the symmetric choice, whose freezing error is
$O(\Delta\xi^2)$---and invert the change of variables at the new column
(the constant envelope $e^{-i\beta\Delta\xi}$ of \eqref{eq:expsymbol}
multiplies every step identically and is carried by the reduced field):
\[
\varphi_{n+1}=\Op\bigl(b_{n+1}^{-1/4}\bigr)\,
e^{i\Delta\xi(\sqrt{B_{n+1/2}}-\beta)}\,
\Op\bigl(b_n^{1/4}\bigr)\,\varphi_n+O(\varepsilon^2\Delta\xi).
\]
Composed by the Moyal product,
\[
b_{n+1}^{-1/4}\,\#\,e^{i\Delta\xi(\sqrt{b_{n+1/2}}-\beta)}\,\#\,b_n^{1/4}
=e^{i\Delta\xi(\sqrt{b_{n+1/2}}-\beta)}
\left(\frac{b_n}{b_{n+1}}\right)^{1/4}+O(\varepsilon\,\Delta\xi),
\]
where the corrections are Poisson brackets between neighbouring
columns---the neglected commutators of Section~\ref{sec:factor}, now per
step,
\[
\{b_n,b_{n+1/2}\}=\frac{\Delta\xi}2\,\{b,\pd_\xi b\}+\dots,\qquad
\text{of order } \varepsilon\,\Delta\xi .
\]
The pointwise part is exactly the one-step WKB map
\[
e^{i\Delta\xi(\sqrt{b_{n+1/2}}-\beta)}
\,\exp\Bigl(-\int_{\xi_n}^{\xi_{n+1}}\rho\;d\xi\Bigr)+O(\Delta\xi^2),
\qquad \rho=\frac{\pd_\xi b}{4b}:
\]
the phase from the midpoint column, the amplitude from the change of
variables \eqref{eq:weq}. In operator form the same composition reads
\begin{equation}\label{eq:step}
\varphi_{n+1}=\underbrace{\bigl[B_{n+1}^{-1/4}B_{n+1/2}^{1/4}\bigr]}_{\text{symbol }(b_{n+1/2}/b_{n+1})^{1/4}}
\,e^{i\Delta\xi\sqrt{B_{n+1/2}}}\,
\underbrace{\bigl[B_{n+1/2}^{-1/4}B_n^{1/4}\bigr]}_{\text{symbol }(b_n/b_{n+1/2})^{1/4}}\varphi_n
+O(\varepsilon^2\Delta\xi),
\end{equation}
equal to $B_{n+1}^{-1/4}e^{i\Delta\xi\sqrt{B_{n+1/2}}}B_n^{1/4}\varphi_n$ up
to the same remainder: inserting the middle pair is legitimate at the symbol
level because $b_{n+1/2}^{1/4}\#b_{n+1/2}^{-1/4}=1+O(\varepsilon^2)$, the
first bracket vanishing exactly. Each bracket is the identity wherever its
columns coincide, in particular in the whole exterior, and has the symbol
$1-\tfrac{\Delta\xi}8\,\pd_\xi b/b+O(\Delta\xi^2)$.

\paragraph{Staircase form.} The column-by-column step replaces the symbol
$b(\xi,\cdot)$ by a staircase of column symbols. At an interface between
$\lambda_a=\sqrt{b_a}$ and $\lambda_b=\sqrt{b_a}(1+\delta)$, matching
$\varphi$ and $\varphi_\xi$ gives the amplitude transmission
$t=(b_a/b_b)^{1/4}\tau$ with
$\tau=2\sqrt{\lambda_a\lambda_b}/(\lambda_a+\lambda_b)=1-\delta^2/8+\dots$,
and a reflected amplitude of relative size
$\delta/(2+\delta)\approx\rho\,\Delta\xi$. The step takes
$\tau=1$---exactly what \eqref{eq:step} does---and so suppresses these
spurious interface reflections, which would otherwise accumulate to
$O(\rho^2\ell_\xi\Delta\xi)$ over a transition of length $\ell_\xi$; the
reflections of the continuous symbol cancel by phase
(Section~\ref{sec:backscatter}). The brackets of \eqref{eq:step} are
thus the discrete counterpart of the leading term of the Bremmer series.

\paragraph{Rotated factors.} The factors $(1+L)^{\pm1/4}$ are approximated by
rational functions; their Pad\'e poles lie on the branch cut $L<-1$, inside
the evanescent spectrum, where a pair of factors at slightly different
columns can amplify a spurious mode (a localized $14\%$ error was observed
after $600$\,m). Following \cite{milinazzo1997}, the approximant is
taken of $(1+Y)^{1/4}$ with $Y=(1+L)e^{-i\gamma}-1$, $\gamma=45^\circ$;
rotated off the real axis, the almost-cancelling poles of almost-cancelling
factors can no longer collide.

\subsection{What the factorization drops}\label{sec:backscatter}
The Bremmer series \cite{bremmer,atkinson,dehoop} ranks the dropped terms.
With $Q=B^{1/4}$ and the normalized amplitudes
$u_\pm=\tfrac12(Q\varphi\mp iQ^{-1}\varphi_\xi)$, the exact split is
\[
\pd_\xi\begin{pmatrix}u_+\\u_-\end{pmatrix}
=\begin{pmatrix}i\Lambda+T & \mathcal R\\ \mathcal R & -i\Lambda+T\end{pmatrix}
\begin{pmatrix}u_+\\u_-\end{pmatrix},\quad
T=\tfrac12[Q_\xi,Q^{-1}],\ \ \mathcal R=\tfrac12(Q_\xi Q^{-1}+Q^{-1}Q_\xi),
\]
whose symbols follow from \eqref{eq:moyal}:
\[
\rho=\operatorname{symbol}(\mathcal R)=\frac{\pd_\xi b}{4b}
\bigl(1+O(\varepsilon^2)\bigr),\qquad
\operatorname{symbol}(T)=\frac i{32}\,\frac{\{b,\pd_\xi b\}}{b^2}
+O(\varepsilon^3).
\]
The coupling $\rho$ is $O(\varepsilon)$ and $T$ is a pure phase, $O(\varepsilon^2)$.
The one-way step is the zeroth Bremmer term. Of the first reflection
\[
u_-(\xi)=-\int_\xi^X\rho\,e^{2i\Phi}\,d\xi',\qquad \Phi_\xi=\lambda_0 ,
\]
integration by parts separates a \emph{slaved, local part}
$\rho\,e^{i\Phi}/(2i\lambda_0)$---a forward-travelling dressing whose
back-action produces exactly the $\lambda_2$ of \eqref{eq:lambda2}, hence
$O(\varepsilon^2)$ and includable in a one-way step---from the
\emph{genuine backscatter}, an oscillatory integral with phase $2\Phi$ and no
stationary point: $O(\varepsilon^\infty)$ for a smooth metric, and
$\approx e^{-2kh_\xi\cos\theta\,\delta}$ for a metric analytic in a strip of
width $\delta$. (Numerical check: for $y''+k^2(1+0.3\tanh(x/\ell))y=0$ the
exact reflection and the first Bremmer term agree to $0.6\%$ over
$k\ell=0.5\dots3$, decaying like $e^{-\pi k\ell}$.) Terrain smoothing is
therefore not cosmetic: a kink in the profile makes the metric non-analytic
and the backscatter only algebraic---edge diffraction, which no one-way method
represents.

%=====================================================================
\section{How the rational approximation enters}\label{sec:rational}

\paragraph{The step and its symbol.} On one column the propagator is a
function of the single operator $L=B/\beta^2-1$. The exponential
\eqref{eq:expsymbol} is replaced by a rational function of degree $(p,p+1)$,
\begin{equation}\label{eq:rational}
e^{i\beta\Delta\xi(\sqrt{1+L}-1)}\ \approx\ R(L)
=\prod_{l=1}^{p}\frac{1+a_lL}{1+b_lL},
\end{equation}
each factor being one tridiagonal solve. By \eqref{eq:funcalc} the
approximation error is, to leading order, a \emph{pointwise function on phase
space}:
\begin{equation}\label{eq:E}
E(\xi,\eta,q)=\Bigl|\,R\!\left(\tfrac b{\beta^2}-1\right)
-e^{i\Delta\xi(\sqrt b-\beta)}\Bigr| .
\end{equation}
There is no operator between the approximation and its target: both are
functions of one operator, and the calculus collapses their comparison to
arithmetic on symbols.

\paragraph{The error budget.} What must be small is $E$ on the support of the
field in phase space---its wavefront set, in practice the support of its
Wigner function. By \eqref{eq:angle}, for angles up to the aperture
$\theta_{\max}$ and wavenumbers in $[k_{\min},k_{\max}]$, the symbol $b$ ranges
over
\[
\bigl[k_{\min}^2h_-^2\cos^2\theta_{\max},\ k_{\max}^2h_+^2\bigr],
\qquad h_\pm=\max/\min h_\xi ,
\]
and these are exactly the rescaled bounds with which the grid optimizer is
called \cite{lytaev2020iccsa,lytaev2026nmp}. That $E$ depends on the metric
only through this range is the precise form of ``local equivalence''.

\paragraph{Per-class reference wavenumber.} On one column $b$ takes values
only in that column's $h_\xi$ range; choosing $\beta_k$ per class makes the
Pad\'e approximant cover a narrower interval of $b/\beta_k^2-1$. The envelope
factor $e^{i(\beta_k-\beta)\Delta\xi}$ is a constant phase and does not change
the symbol.

\paragraph{Where the field lives: bicharacteristics.} The support of the field
is not a property of one column; it is transported by the Hamiltonian flow of
the one-way symbol,
\begin{equation}\label{eq:rays}
\frac{d\eta}{d\xi}=-\pd_q\sqrt{b_0},\qquad
\frac{dq}{d\xi}=\pd_\eta\sqrt{b_0},
\end{equation}
and reflected at the ground by the boundary symbol (the discrete Robin
condition of Section~\ref{sec:scheme}). A ray reflected by a slope
$\theta_s$ leaves with up to $\theta_{\max}+2\theta_s$ to the horizontal, and
once the terrain flattens, the local symbol range is that of a flat column
while the field still carries the large $q$ produced upstream. Error budgets
must therefore be evaluated on the \emph{transported} support---the
slope-widened aperture. This is not academic: per-class range steps sized for
their own columns' symbol ranges failed until the aperture was widened, and
the wide-angle floors of the shift map (Section~\ref{sec:obstructions}) are
the same statement in reverse.

\paragraph{Contrast with the classical schemes.} The Crank--Nicolson
discretization of shift-map PEs first approximates the square root,
$u_x=ik_0(f(X)-1)u$ with $f\approx\sqrt{1+X}$, $X=k_0^{-2}\pd_\zeta^2+\dots$,
and then the exponential by the Cayley transform. Its error thus contains
both a square-root error and a $\Delta x$ error, growing with angle
independently of the order of $f$. The rational step \eqref{eq:rational}
approximates the propagator itself, to order $2p$ in $L$ on the band that
matters, and its error vanishes uniformly as $p$ grows for range steps of
many wavelengths.

%=====================================================================
\section{The numerical scheme}\label{sec:scheme}

\paragraph{Vertical discretization.} On the uniform grid $\eta_j=j\Delta\eta$
each factor of \eqref{eq:step}, with $\varphi=r^{1/2}\psi$,
$m=r^2$, $\het=(k^2h_\xi^2-V)/\beta^2-1$ and the operator form of
Section~\ref{sec:symbol}, is the weighted scheme
\begin{multline}\label{eq:scheme}
c\,\delta^2_rv+(1+\alpha\delta^2_r)\bigl[\tau^2m(1+c\het)v\bigr]\\
=a\,\delta^2_ru+(1+\alpha\delta^2_r)\bigl[\tau^2m(1+a\het)u\bigr]
\end{multline}
with $\tau=\beta\Delta\eta$ and the conservatively discretized weighted
Laplacian
\[
\delta^2_r u_j=r_j\left[\frac{u_{j+1}-u_j}{r_{j+1/2}}-\frac{u_j-u_{j-1}}{r_{j-1/2}}\right] .
\]
For $g=\ln r$ and $D=r\pd r^{-1}\pd$, Taylor expansion gives
$\delta^2_r/\Delta\eta^2=D+\tfrac{\Delta\eta^2}{12}D^2+\Delta\eta^2E+O(\Delta\eta^4)$.
The Numerov weight $\alpha=\tfrac1{12}$ cancels the $D^2$ term, so the scheme
is fourth-order in the wave part, with a second-order remainder of relative
size $(\Delta\eta/\ell_\eta)^2/24$, negligible since $\ell_\eta\gtrsim L$.
This stencil is the discrete form of the operator in \eqref{eq:liou}: it
carries the zeroth-order symbol correction $W=r_\eta^2/(4r^4)$ of
\eqref{eq:symbol} implicitly, exactly as its continuous counterpart does.

\paragraph{Ground closure.} The Dirichlet unknown is eliminated from the
system rather than imposed by a unit row (a unit row next to entries of order
$\alpha\tau^2$ is satisfied by a pivoting solver only up to amplified
round-off); Neumann uses a mirror node. For the Robin condition a ghost node
$u_{-1}$ is closed by the compact (fourth-order) relation
\begin{equation}\label{eq:closure}
(u_1-u_{-1})+\frac g3(u_{-1}+4u_0+u_1)=0,\qquad g=\gamma\Delta\eta ,
\end{equation}
which combines the compact approximations of $\psi'$ and $\psi''$
\cite{lele1992compact}; frozen coefficients annihilate
$\pd_\eta^2(\psi'+\gamma\psi)$, making the closure fourth-order accurate. It
expresses $u_{-1}$ through $u_0,u_1$, so the system stays tridiagonal, and
the same closure serves all $p+4q$ factors of a step.

\paragraph{Top boundary.} Every factor is closed at the top by the discrete
NLBC \cite{lytaev2018awpl},
\begin{equation}\label{eq:nlbc}
u^n_{N-1,l}-\sum_{l'\le l}Q^{(0)}_{ll'}u^n_{N,l'}
=\sum_{m=1}^{n}\sum_{l'}Q^{(m)}_{ll'}u^{n-m}_{N,l'} ,
\end{equation}
whose kernels are the inverse $z$-transform of a matrix function of the
rational coefficients and of the discrete exterior equation; they depend only
on the exterior scheme and on the chain of factors. Its Weyl-symbol reading
is the subject of Section~\ref{sec:top}.

\paragraph{One range step.} The coefficients $\het$, $m$, $r$ are tabulated on
coarse columns below the Cartesian region and interpolated linearly in range;
the main chain is evaluated at the midpoint.
\begin{center}
\begin{minipage}{0.92\linewidth}
\hrule\vspace{2pt}
\textbf{Algorithm 1:} one range step of \eqref{eq:step}--\eqref{eq:scheme}
\vspace{1pt}\hrule\vspace{2pt}
\textbf{Input:} $\varphi_n$ (equivalently $\psi_n$), columns $\het,m,r$ at
$\xi_n,\xi_{n+1/2},\xi_{n+1}$
\begin{enumerate}
\item apply the first conjugating bracket $C_1=B_{n+1/2}^{-1/4}B_n^{1/4}$
  \emph{(rows below $\etal$ plus a margin only)};
\item for $l=1,\dots,p$: solve \eqref{eq:scheme} with $(a_l,b_l)$ at
  $\xi_{n+1/2}$ and the NLBC row \eqref{eq:nlbc};
\item apply the second conjugating bracket $C_2=B_{n+1}^{-1/4}B_{n+1/2}^{1/4}$.
\end{enumerate}
\vspace{-4pt}\hrule
\end{minipage}
\end{center}
One step costs $p+4q$ tridiagonal solves, the $4q$ weight solves restricted
to the transition layer; both brackets are the identity in the exterior,
where the step degenerates to the flat scheme \cite{lytaev2020awpl}.

%=====================================================================
\section{Dispersion relations as Weyl symbols}\label{sec:dispersion}

The scheme has two dispersion relations, and both are symbols.

\paragraph{Continuous.} The principal symbol \eqref{eq:b0} of $B$ is the
continuous dispersion relation, $\kappa^2=k^2h_\xi^2-q^2/r^2$; in angle form,
$\kappa=\pm kh_\xi\cos\theta$. It tells the propagator what $\sqrt b$ to
approximate and the budget what band $\bigl[k_{\min}^2h_-^2\cos^2\theta_{\max},
k_{\max}^2h_+^2\bigr]$ to cover.

\paragraph{Discrete.} Freezing the coefficients and substituting
$\psi^n_j=\lambda^ne^{iqj\Delta\eta}$ into \eqref{eq:scheme} gives the
\emph{discrete Weyl symbol} on the Brillouin band
$q\in(-\pi/\Delta\eta,\pi/\Delta\eta]$,
\begin{equation}\label{eq:Lh}
L_h(q)=\het-\frac4{\tau^2r^2}\,\frac{s}{1-4\alpha s},\qquad
s=\sin^2\frac{q\Delta\eta}2,\qquad
\lambda=R\bigl(L_h(q)\bigr),
\end{equation}
a periodic function agreeing with $b/\beta^2-1=\het-q^2/(\beta r)^2$ to
$O((q\Delta\eta)^6)$ for $\alpha=1/12$. The step's symbol is $R(L_h(q))$, and
the dispersion error per unit physical length is
\begin{equation}\label{eq:disp}
\varepsilon(\theta)=\frac{\bigl|\log\bigl(\lambda\,e^{-i\Delta\xi(p-\beta)}\bigr)\bigr|}
{h_\xi\Delta\xi},\qquad p=h_\xi k\cos\theta ,
\end{equation}
the distance between the discrete symbol and the exact one \eqref{eq:expsymbol}.

\paragraph{Local equivalence.} In physical variables
$\beta_e=\beta/h_\xi$, $\Delta x_e=h_\xi\Delta\xi$, $\Delta z_e=h_\eta\Delta\eta$
one has $\het=k^2/\beta_e^2-1$, $\tau r=\beta_e\Delta z_e$,
$q\Delta\eta=k_z\Delta z_e$, $\beta\Delta\xi=\beta_e\Delta x_e$, hence
\begin{equation}\label{eq:equiv}
\varepsilon(\theta;\beta,\Delta\xi,\Delta\eta,h_\xi,h_\eta)
=\varepsilon_0(\theta;\beta_e,\Delta x_e,\Delta z_e),
\end{equation}
the dispersion error of the \emph{flat} scheme on the effective grid: the
discrete symbol of a distorted cell is the flat discrete symbol. (Verified
numerically to $4\cdot10^{-13}$ of the maximum error.) Three consequences.

\begin{itemize}
\item \emph{No angle restriction from the transformation.} A wave is resolved
  iff its angle to the \emph{local} marching direction lies in the aperture
  of $R$ on the local grid. Waves reflected by slopes up to $\theta_s$ reach
  $\theta_{\max}+2\theta_s$ (Section~\ref{sec:rational}), so the aperture is
  set to that value---a physical requirement that a staircase faces equally.
\item \emph{Grid selection on flat terrain.} The flat-terrain optimizer
  \cite{lytaev2020iccsa,lytaev2026nmp} is called with the wavenumber interval
  $[k_{\min}h_-,k_{\max}h_+]$, the vertical-wavenumber bound $k_z^{\max}h_-$
  and the error budget per unit computational length; the vertical step it
  returns is divided by $\max r\cdot h_+/h_-$. Equivalently, the stretching
  rule $\max_\xi h_\eta s\,\Delta\eta\le\Delta z_{\rm phys}$: the error of
  $L_h$ is a function of $q\Delta\eta=k\sin\theta\cdot h_\eta s\,\Delta\eta$,
  so it is the \emph{physical} step that must stay bounded at every level.
\item \emph{Per-class steps.} Since the symbol depends on the metric only
  through the products $h_\xi\Delta\xi$, $h_\eta\Delta\eta$, $\beta/h_\xi$,
  per-class $\beta_k$, Pad\'e chains and range steps are admissible---but the
  error budget must be evaluated on the transported support of
  \eqref{eq:rays}, not on each class's own symbol range.
\end{itemize}

%=====================================================================
\section{Two obstructions}\label{sec:obstructions}

The same calculus explains two failures that no amount of grid refinement
cures: the error floor of the shift map, and the inability of any local top
boundary to absorb grazing waves. Both are mismatches of \emph{symbols}, hence
independent of the discretization.

\subsection{Why the shift map has a grid-independent floor}\label{sec:shift}

Write the shift-map equation on a constant slope $s=T'$ as
\begin{equation}\label{eq:shiftmap}
\psi_{xx}-2s\,\psi_{x\zeta}+(1+s^2)\,\psi_{\zeta\zeta}+k^2\psi=0,
\qquad \zeta=z-T(x),
\end{equation}
whose symbol in the duals $(\kappa,q)$ of $(x,\zeta)$ is
$-\kappa^2+2s\kappa q-(1+s^2)q^2+k^2$. The outgoing root is
\begin{equation}\label{eq:shiftroot}
\kappa=sq+\sqrt{k^2-q^2} .
\end{equation}

\paragraph{The obstruction.} The one-way symbol has a part $sq$ \emph{odd in
$q$}. But the exact one-way symbol in orthogonal coordinates is
$\sqrt{b_0}$---a function of the symbol of the single symmetric second-order
vertical operator $B$, which has no first-order part---and no function of a
single symmetric vertical operator has a symbol with an odd part. The shift
map therefore cannot be represented by a propagator of the form
$F(\text{vertical operator})$ at any order: the deficiency is structural, not
a matter of approximation order.

\paragraph{The gauge.} The Beilis--Tappert gauge
\cite{beilis1979coupled} $u=v\,e^{ik_0T'\zeta+ik_0\int T'^2dx'/2}$ translates
the dual variable, $q=q'+k_0s$, and removes the range phase; the transformed
root becomes
\[
\kappa'=sq'+\frac{k_0s^2}2+\sqrt{k^2-(q'+k_0s)^2}-k_0 .
\]
If the square root is truncated at second order (the narrow-angle symbol), the
terms linear in $q'$ and the constants cancel and
$\kappa'=-q'^2/(2k_0)$ \emph{exactly}: the gauge is exact for the narrow-angle
PE. The higher terms involve $(q'+k_0s)^4,\dots$ and do not cancel.

\paragraph{The floor.} A wide-angle PE propagates with the symbol
$k_0[f(-q'^2/k_0^2)-1]$, an \emph{even} function of $q'$, while the exact
$\kappa'$ is not even beyond second order. The difference is a mismatch of
symbols of order $O(1)$ in $\varepsilon$---independent of $\Delta x$,
$\Delta\zeta$, of the stencil order and of the accuracy of $f$. In terms of
the angle $\alpha$ to the horizontal on the slope, after removing the
angle-independent part $f(0)-1$,
\begin{equation}\label{eq:floor}
\varepsilon_\infty(\alpha;\theta_s)
=k_0\Bigl|f\bigl(-(\sin\alpha-s)^2\bigr)+\tfrac{s^2}2-s\sin\alpha-\cos\alpha\Bigr| .
\end{equation}
For $f=1+X/2$ this reduces to the flat-ground NA floor,
$k_0|1-\tfrac12\sin^2\alpha-\cos\alpha|$: the BT transformation is exact for
the NA PE, and NA on a slope is NA on flat ground. For every other $f$,
\emph{including the exact square root}, the floor depends on the slope and
vanishes only on a narrow set of angles. At a budget of $10^{-4}$\,rad$/\lambda$
the exact square root admits only waves between $3.9^\circ$ and $6.3^\circ$
above the horizontal on a $10^\circ$ slope, a single angle on $20^\circ$, and
none on $30^\circ$; raising the order of $f$ cannot help, because the floor
is set by the gauge. This quantifies the published observations that WA
shift-map PEs are limited to slopes of $10$--$20^\circ$ and can be worse than
the NA PE on lee-side slopes \cite{wang2020smpe,zhou2017lfhighangle}, and that
their sheared vertical wavenumber is unphysical \cite{gilbert2022wabt}. In the
orthogonal coordinates of this paper the frozen-coefficient error is, by
\eqref{eq:equiv}, that of the flat scheme on the local effective grid, and it
vanishes under refinement for every slope---because the exact one-way symbol
$\sqrt{b_0}$ is a function of the symbol of one operator, the very property
whose absence created the floor.

\subsection{Why no local top boundary absorbs grazing waves}\label{sec:top}

\paragraph{Continuous exterior.} Above the transition layer the medium is
invariant in $x$, and the field there is a superposition of
$e^{i\kappa x}e^{i\sqrt{k^2-\kappa^2}\,z}$; the transparent condition is
\begin{equation}\label{eq:tnlbc}
\pd_z\psi=i\,\Op_x\!\left(\sqrt{k^2-\kappa^2}\right)\psi .
\end{equation}
This is a pseudodifferential operator in the \emph{evolution} variable, hence
a convolution in range: the exact condition carries the whole history of the
boundary values. That is the definition of the NLBC \eqref{eq:nlbc}. Its
symbol has branch points at $\kappa=\pm k$, i.e.\ at vertical wavenumber
$k_z=0$: waves that graze the top.

\paragraph{Discrete exterior.} In the scheme, the $z$-transform variable $t$
plays the role of $e^{i\kappa\Delta x}$, and the scalar exterior ratio
$R(\lambda)=u_{N-1}/u_N$, with $\lambda=\beta^2(1+\sigma)$, is the discrete
symbol of the condition. It has a discrete branch point where the two roots
$\mu=e^{\pm ik_z\Delta z}$ meet---again at $k_z=0$.

\paragraph{Local conditions are rational approximations.} A boundary condition
that is local in range---an absorbing layer, a complex-index or complex-step
layer of $M$ rows, a series of local ABCs---has an exterior ratio that is a
\emph{rational function} of $\lambda$ with finitely many poles. A rational
function cannot approximate a square root uniformly up to its branch point
\cite{newman,idk}:
\begin{itemize}
\item on the band $k_z\Delta z\in[\theta_{\min},\theta_{\max}]$ the best
  error decays geometrically in $M$, at a rate set by
  $\log(\theta_{\max}/\theta_{\min})$;
\item as $\theta_{\min}\to0$ the rate degenerates to zero.
\end{itemize}
The continuous picture says the same: a layer of complex thickness $\Lambda$
reflects about $e^{-2k_z\operatorname{Im}\Lambda}$, which tends to $1$ as
$k_z\to0$ for any finite layer. Grazing waves feel no layer that ends. The
complex layer implemented in the code is the optimal-grid construction of
\cite{idk,adgk} applied to the Numerov exterior, and its row counts
($8$, $12$, $16$ and $34$ rows for $5\cdot10^{-3}$, $2\cdot10^{-4}$,
$10^{-5}$, $10^{-6}$ over $0.5$--$30^\circ$) show exactly this geometric rate.
Numerically, with everything else fixed, the NLBC leaves errors of
$4\cdot10^{-7}$--$4.6\cdot10^{-3}$ for beams at $1$--$30^\circ$ (the last
being the grid's propagation error, shared by every boundary), while the best
absorbing-layer settings give $1.6\cdot10^{-3}$ at $1^\circ$---four thousand
times worse---and no setting serves all angles, because each angle band needs
its own rational approximation.

\paragraph{Why one condition serves every class.} For per-class $\beta_k$,
Pad\'e chains and range steps, the exterior equation of every substep is
$(A-\lambda)w=0$ with the \emph{same} $A$. A condition expressed as a function
of $\lambda$---the symbol of $A$ alone---is therefore valid for all of them.
The convolution NLBC is instead a function of $t$, and the relation between
$t$ and $\lambda$ changes with the class; it must then be computed per class
or replaced by the $\lambda$-form. This is the boundary-condition counterpart
of the per-class budget of Section~\ref{sec:rational}.

\paragraph{Summary.} Absorbing a grazing wave transparently requires either a
condition with an infinite memory (the NLBC) or an infinite stretch. No
finite, local, range-independent condition can do it, and the obstruction is
visible in the symbol: a branch point that every rational approximation must
keep at a finite distance.

%=====================================================================
\section{Summary of orders}\label{sec:orders}

\begin{center}
\footnotesize
\begin{tabular}{@{}lll@{}}
\toprule
term & symbol & order \\
\midrule
dispersion relation & $b_0=k^2h_\xi^2-q^2/r^2$ & $1$ \\
first-order (transport) part of $B$ & $0$ (Liouville gauge) & --- \\
$V$, $W$ & $(r^{1/2})_{\xi\xi}/r^{1/2}$,\ \ $r_\eta^2/(4r^4)$ & $\varepsilon^2$ \\
transport term of the one-way eq. & $\tfrac i4 b_\xi/b$ & $\varepsilon$ (removed by $B^{1/4}$) \\
second-order WKB, commutators & \eqref{eq:lambda2}, $i\{b,b_\xi\}$ & $\varepsilon^2$ \\
conjugating brackets at neighbouring columns & $\{b_n,b_{n+1/2}\}$ & $\varepsilon\,\Delta\xi$ \\
Bremmer coupling $\mathcal R$ & $b_\xi/(4b)$ & $\varepsilon$ \\
\quad slaved part (local) & $-\rho^2/(2\lambda_0)$, part of $\lambda_2$ & $\varepsilon^2$ \\
\quad genuine backscatter & $\int\rho\,e^{2i\Phi}d\xi$ & $\varepsilon^\infty$ ($e^{-2kh\cos\theta\,\delta}$) \\
staircase reflection (avoided) & $\rho^2\Delta\xi^2$ per step & $\varepsilon^2\Delta\xi$ total \\
rational step vs exponential & $R(b/\beta^2-1)-e^{i\Delta\xi(\sqrt b-\beta)}$ & set by the optimizer \\
Numerov symbol error & $L_h(q)-(b/\beta^2-1)$ & $(q\Delta\eta)^6$ \\
shift map, wide angle & odd part of $sq+\sqrt{k^2-q^2}$ & $O(1)$, grid independent \\
local top condition & rational approx.\ of $\sqrt{\cdot}$ near $k_z=0$ & geometric in $M$, not uniform \\
\bottomrule
\end{tabular}
\end{center}

Every line is a statement about one function on phase space. The method is
exact where its symbols agree with the exact ones, and its two irreducible
limitations---the turning set $b_0=0$ that bounds the aperture, and the
branch point $k_z=0$ that forces a nonlocal top boundary---are visible in the
same language as the scheme itself.

\begin{thebibliography}{24}

\bibitem{levy2000book}
M.~Levy, \emph{Parabolic Equation Methods for Electromagnetic Wave
  Propagation}, ser. IEE Electromagnetic Waves Series 45. London: IEE, 2000.

\bibitem{lytaev2020awpl}
M.~S. Lytaev, ``{Numerov--Pad\'e} scheme for the one-way {Helmholtz} equation
  in tropospheric radio-wave propagation,'' \emph{IEEE Antennas and Wireless
  Propagation Letters}, vol.~19, no.~12, pp. 2167--2171, 2020.

\bibitem{lytaev2018awpl}
------, ``Nonlocal boundary conditions for split-step {Pad\'e} approximations
  of the {Helmholtz} equation with modified refractive index,'' \emph{IEEE
  Antennas and Wireless Propagation Letters}, vol.~17, no.~8, pp. 1561--1565,
  2018.

\bibitem{lytaev2020iccsa}
------, ``Automated selection of the computational parameters for the
  higher-order parabolic equation numerical methods,'' in \emph{Computational
  Science and Its Applications -- ICCSA 2020}, ser. LNCS, vol. 12249.
  Springer, 2020, pp. 296--311.

\bibitem{lytaev2026nmp}
------, ``Mesh optimization for rational approximations of the one-way
  {Helmholtz} equation in inhomogeneous media,'' \emph{Numerical Methods and
  Programming (Vychislitel'nye Metody i Programmirovanie)}, vol.~27, no.~2,
  pp. 302--316, 2026.

\bibitem{beilis1979coupled}
A.~Beilis and F.~D. Tappert, ``Coupled mode analysis of multiple rough surface
  scattering,'' \emph{J. Acoust. Soc. Am.}, vol.~66, no.~3, pp. 811--826, 1979.

\bibitem{barrios1994tpem}
A.~E. Barrios, ``A terrain parabolic equation model for propagation in the
  troposphere,'' \emph{IEEE Trans. Antennas Propag.}, vol.~42, no.~1,
  pp. 90--98, 1994.

\bibitem{donohue2000gtpe}
D.~J. Donohue and J.~R. Kuttler, ``Propagation modeling over terrain using the
  parabolic wave equation,'' \emph{IEEE Trans. Antennas Propag.}, vol.~48,
  no.~2, pp. 260--277, 2000.

\bibitem{holm2007wasm}
P.~D. Holm, ``Wide-angle shift-map {PE} for a piecewise linear terrain---a
  finite-difference approach,'' \emph{IEEE Trans. Antennas Propag.}, vol.~55,
  no.~10, pp. 2773--2789, 2007.

\bibitem{wang2020smpe}
D.-D. Wang, Y.-R. Pu, X.-L. Xi, and L.-L. Zhou, ``An analysis of narrow-angle
  and wide-angle shift-map parabolic equations,'' \emph{IEEE Trans. Antennas
  Propag.}, vol.~68, no.~5, pp. 3911--3918, 2020.

\bibitem{zhou2017lfhighangle}
L.~Zhou, D.~Wang, Z.~Mu, X.~Xi, and L.~He, ``{LF} radio wave prediction at
  short ranges with high propagation angles over irregular terrain,''
  \emph{IEEE Antennas and Wireless Propagation Letters}, vol.~16,
  pp. 732--735, 2017.

\bibitem{gilbert2022wabt}
K.~E. Gilbert, ``Wide-angle formulation of the {Beilis--Tappert} method,''
  \emph{J. Acoust. Soc. Am.}, vol. 152, no.~2, pp. 1170--1179, 2022.

\bibitem{milinazzo1997}
F.~A. Milinazzo, C.~A. Zala, and G.~H. Brooke, ``Rational square-root
  approximations for parabolic equation algorithms,'' \emph{J. Acoust. Soc.
  Am.}, vol. 101, no.~2, pp. 760--766, 1997.

\bibitem{lele1992compact}
S.~K. Lele, ``Compact finite difference schemes with spectral-like
  resolution,'' \emph{J. Comput. Phys.}, vol. 103, no.~1, pp. 16--42, 1992.

\bibitem{zworski}
M.~Zworski, \emph{Semiclassical Analysis}. Providence, RI: AMS, 2012.

\bibitem{dimassi}
M.~Dimassi and J.~Sj\"ostrand, \emph{Spectral Asymptotics in the
  Semi-Classical Limit}. Cambridge Univ. Press, 1999.

\bibitem{helffer}
B.~Helffer and D.~Robert, ``Calcul fonctionnel par la transformation de
  Mellin et op\'erateurs admissibles,'' \emph{J. Funct. Anal.}, vol.~53,
  pp. 246--268, 1983.

\bibitem{fishman}
L.~Fishman and J.~J. McCoy, ``Derivation and application of extended parabolic
  wave theories,'' \emph{J. Math. Phys.}, vol.~25, pp. 285--308, 1984.

\bibitem{dehoop}
M.~V. de~Hoop, ``Generalization of the Bremmer coupling series,'' \emph{J.
  Math. Phys.}, vol.~37, pp. 3246--3282, 1996.

\bibitem{bremmer}
H.~Bremmer, ``The W.K.B. approximation as the first term of a geometric-optical
  series,'' \emph{Comm. Pure Appl. Math.}, vol.~4, pp. 105--115, 1951.

\bibitem{atkinson}
F.~V. Atkinson, ``Wave propagation and the Bremmer series,'' \emph{J. Math.
  Anal. Appl.}, vol.~1, pp. 255--276, 1960.

\bibitem{newman}
D.~J. Newman, ``Rational approximation to $|x|$,'' \emph{Michigan Math. J.},
  vol.~11, pp. 11--14, 1964.

\bibitem{idk}
D.~Ingerman, V.~Druskin and L.~Knizhnerman, ``Optimal finite difference grids
  and rational approximations of the square root. I. Elliptic problems,''
  \emph{Comm. Pure Appl. Math.}, vol.~53, pp. 1039--1066, 2000.

\bibitem{adgk}
S.~Asvadurov, V.~Druskin, M.~N. Guddati and L.~Knizhnerman, ``On optimal
  finite-difference approximation of {PML},'' \emph{SIAM J. Numer. Anal.},
  vol.~41, pp. 287--305, 2003.

\end{thebibliography}

\end{document}
